\scanpage{12}
\mspage{9}
\uline{et pour différentielle invariante} $\pi=dx/2y$.

De plus, sous ces hypothèses, la donnée d'une forme différentielle invariante suffit à rigidifier une courbe.

4.\qquad \uline{Cas où} 2 \uline{est nilpotent et} $c_4$ \uline{inversible}.

Lorsque $2=0$, la forme différentielle ne suffit plus jamais à rigidifier une courbe de genre un. Lorsque $c_4$ est inversible et 2 nilpotent, donc $a_1$ inversible, en changeant le choix de $\pi$, $x$, $y$ on peut trouver une équation
\[
y^2+xy=x^3+a_2x^2+a_6\qquad \pi=\frac{dx}{x+2y}\qquad .\tag*{(4.1)}
\]
Les autres choix de $x$ et $y$ donnant lieu à une telle équation s'obtiennent en faisant, dans (1.2)
\[
r=t=0\qquad u=1+2s\tag*{(4.2)}
\]
et on aura
\[
\begin{aligned}
u^2a'_2&=a_2-s-s^2\\
u^6a'_6&=a_6
\end{aligned}\tag*{(4.3)}
\]
et \uline{modulo 2}
\[
c_4=1\qquad c_6=1\qquad \Delta=a_6\qquad .\tag*{(4.4)}
\]

5.\qquad \uline{Cas où} $c_4=0$, $p=2$ \uline{ou} $p=3$

a) $p=3$

En carac\deleted{\unsure{e|b|t}}téristique 3 , lorsque l'équation est mise sour forme

% archive sheet number: 11
